% GENERATED FILE. Edit canonical lessons, then run npm run generate:books.
% canonical-source-hash: fnv1a64:362dca37e6615a40
% canonical-lessons: UR-C412-qasba, UR-C412-abadi, UR-C412-bashinda, UR-C412-dehat, UR-C412-paros

\chapter{Towns and Villages}
\label{ch:ur-towns-and-villages}
\hlchaptermodality{412}

\begin{chapteropening}
\textbf{By the end of this chapter:} \emph{I can say where I live: a small town, the countryside, the people next door.}

Complete the last lesson of chapter 412: I can say where I live: a small town, the countryside, the people next door.
\end{chapteropening}

\section[\texorpdfstring{\ur{قصبہ} (qasbā)}{قصبہ (qasbā)}]{\ur{قصبہ} (qasbā) — a small town}

\label{lesson:UR-C412-qasba}



\begin{quote}
Before the new one: say the Urdu for a basement, then the Urdu for beside.
\end{quote}

\subsection*{You'll want to know: \ur{قصبہ}}
\textbf{\ur{قصبہ}} — \emph{qasbā} — \textquotedblleft{}a small town\textquotedblright{}.

\begin{quote}
\ur{ہمارا} \ur{قصبہ} \ur{شہر} \ur{سے} \ur{دور} \ur{ہے۔}
\end{quote}

\begin{quote}
\emph{hamārā qasbā shahr se dūr hai.}
\end{quote}

\begin{quote}
Our little town is a long way from the city.
\end{quote}

Masculine.

\begin{cousinweb}
\textbf{\ur{قصبہ}} (\emph{qasbā}) is usually traced, through Persian, to Arabic \emph{qaṣaba}, \textquotedblleft{}a chief town, a citadel\textquotedblright{}; in Urdu it is a town smaller than a \textbf{\ur{شہر}}.
\end{cousinweb}

\begin{grammarlens}[title={What you've built}]
One more word for getting around.
\end{grammarlens}

\subsection*{Your turn}
\begin{itemize}
\raggedright
  \item \emph{Say it:} \emph{qasbā}
  \item \emph{Say it:} \emph{qasbā}, once more
  \item \emph{Say it:} \emph{hamārā qasbā shahr se dūr hai.} — Our little town is a long way from the city.
  \item \emph{Recall:} say the Urdu for a basement, then the Urdu for beside, then say \emph{qasbā} again
\end{itemize}

\begin{tcolorbox}[breakable,colback=teal!4,colframe=teal!35!black,title={Before you move on}]
What does \ur{قصبہ} mean? (\textquotedblleft{}A small town\textquotedblright{}.) Say it once more.
\end{tcolorbox}
\section[\texorpdfstring{\ur{آبادی} (ābādī)}{آبادی (ābādī)}]{\ur{آبادی} (ābādī) — a population; a settlement}

\label{lesson:UR-C412-abadi}



\begin{quote}
Before the new one: say the Urdu for beside, then the Urdu for a small town.
\end{quote}

\subsection*{You'll want to know: \ur{آبادی}}
\textbf{\ur{آبادی}} — \emph{ābādī} — \textquotedblleft{}a population; a settlement\textquotedblright{}.

\begin{quote}
\ur{اس} \ur{قصبے} \ur{کی} \ur{آبادی} \ur{کم} \ur{ہے۔}
\end{quote}

\begin{quote}
\emph{is qasbe kī ābādī kam hai.}
\end{quote}

\begin{quote}
Not many people live in this little town.
\end{quote}

Feminine.

\begin{cousinweb}
\textbf{\ur{آبادی}} (\emph{ābādī}) is built on \emph{ābād}, \textquotedblleft{}inhabited, flourishing\textquotedblright{}, usually traced to Persian \emph{ābād}, with the ending \emph{-ī}.
\end{cousinweb}

\begin{grammarlens}[title={What you've built}]
One more word for getting around.
\end{grammarlens}

\subsection*{Your turn}
\begin{itemize}
\raggedright
  \item \emph{Say it:} \emph{ābādī}
  \item \emph{Say it:} \emph{ābādī}, once more
  \item \emph{Say it:} \emph{is qasbe kī ābādī kam hai.} — Not many people live in this little town.
  \item \emph{Recall:} say the Urdu for beside, then the Urdu for a small town, then say \emph{ābādī} again
\end{itemize}

\begin{tcolorbox}[breakable,colback=teal!4,colframe=teal!35!black,title={Before you move on}]
What does \ur{آبادی} mean? (\textquotedblleft{}A population; a settlement\textquotedblright{}.) Say it once more.
\end{tcolorbox}
\section[\texorpdfstring{\ur{باشندہ} (bāshindā)}{باشندہ (bāshindā)}]{\ur{باشندہ} (bāshindā) — an inhabitant; a resident}

\label{lesson:UR-C412-bashinda}



\begin{quote}
Before the new one: say the Urdu for a small town, then the Urdu for a population.
\end{quote}

\subsection*{You'll want to know: \ur{باشندہ}}
\textbf{\ur{باشندہ}} — \emph{bāshindā} — \textquotedblleft{}an inhabitant; a resident\textquotedblright{}.

\begin{quote}
\ur{وہ} \ur{اسی} \ur{قصبے} \ur{کا} \ur{باشندہ} \ur{ہے۔}
\end{quote}

\begin{quote}
\emph{voh usī qasbe kā bāshindā hai.}
\end{quote}

\begin{quote}
He comes from that very town.
\end{quote}

Masculine or feminine, by the person meant.

\begin{cousinweb}
\textbf{\ur{باشندہ}} (\emph{bāshindā}) is usually traced to Persian \emph{bāshinda}, \textquotedblleft{}one who lives somewhere\textquotedblright{}.
\end{cousinweb}

\begin{grammarlens}[title={What you've built}]
One more word for getting around.
\end{grammarlens}

\subsection*{Your turn}
\begin{itemize}
\raggedright
  \item \emph{Say it:} \emph{bāshindā}
  \item \emph{Say it:} \emph{bāshindā}, once more
  \item \emph{Say it:} \emph{voh usī qasbe kā bāshindā hai.} — He comes from that very town.
  \item \emph{Recall:} say the Urdu for a small town, then the Urdu for a population, then say \emph{bāshindā} again
\end{itemize}

\begin{tcolorbox}[breakable,colback=teal!4,colframe=teal!35!black,title={Before you move on}]
What does \ur{باشندہ} mean? (\textquotedblleft{}An inhabitant; a resident\textquotedblright{}.) Say it once more.
\end{tcolorbox}
\section[\texorpdfstring{\ur{دیہات} (dehāt)}{دیہات (dehāt)}]{\ur{دیہات} (dehāt) — the countryside; villages}

\label{lesson:UR-C412-dehat}



\begin{quote}
Before the new one: say the Urdu for a population, then the Urdu for an inhabitant.
\end{quote}

\subsection*{You'll want to know: \ur{دیہات}}
\textbf{\ur{دیہات}} — \emph{dehāt} — \textquotedblleft{}the countryside; villages\textquotedblright{}.

\begin{quote}
\ur{دیہات} \ur{میں} \ur{لوگ} \ur{کھیتوں} \ur{میں} \ur{کام} \ur{کرتے} \ur{ہیں۔}
\end{quote}

\begin{quote}
\emph{dehāt meṅ log khetoṅ meṅ kām karte haiṅ.}
\end{quote}

\begin{quote}
In the countryside people work in the fields.
\end{quote}

Masculine, used in the plural.

\begin{cousinweb}
\textbf{\ur{دیہات}} (\emph{dehāt}) is built on Persian \emph{deh}, \textquotedblleft{}a village\textquotedblright{}, with the Arabic plural ending \emph{-āt}: it is usually traced to Persian \emph{dehāt}, \textquotedblleft{}villages\textquotedblright{}.
\end{cousinweb}

\begin{grammarlens}[title={What you've built}]
One more word for getting around.
\end{grammarlens}

\subsection*{Your turn}
\begin{itemize}
\raggedright
  \item \emph{Say it:} \emph{dehāt}
  \item \emph{Say it:} \emph{dehāt}, once more
  \item \emph{Say it:} \emph{dehāt meṅ log khetoṅ meṅ kām karte haiṅ.} — In the countryside people work in the fields.
  \item \emph{Recall:} say the Urdu for a population, then the Urdu for an inhabitant, then say \emph{dehāt} again
\end{itemize}

\begin{tcolorbox}[breakable,colback=teal!4,colframe=teal!35!black,title={Before you move on}]
What does \ur{دیہات} mean? (\textquotedblleft{}The countryside; villages\textquotedblright{}.) Say it once more.
\end{tcolorbox}
\section[\texorpdfstring{\ur{پڑوس} (paṛos)}{پڑوس (paṛos)}]{\ur{پڑوس} (paṛos) — next door; the neighbourhood}

\label{lesson:UR-C412-paros}



\begin{quote}
Before the new one: say the Urdu for an inhabitant, then the Urdu for the countryside.
\end{quote}

\subsection*{You'll want to know: \ur{پڑوس}}
\textbf{\ur{پڑوس}} — \emph{paṛos} — \textquotedblleft{}next door; the neighbourhood\textquotedblright{}.

\begin{quote}
\ur{ہمارے} \ur{پڑوس} \ur{میں} \ur{ایک} \ur{درزی} \ur{رہتا} \ur{ہے۔}
\end{quote}

\begin{quote}
\emph{hamāre paṛos meṅ ek darzī rahtā hai.}
\end{quote}

\begin{quote}
A tailor lives next door to us.
\end{quote}

Masculine.

\begin{cousinweb}
\textbf{\ur{پڑوس}} (\emph{paṛos}) is inherited: CDIAL lists the Hindi cousin \emph{paṛos}, \textquotedblleft{}the neighbourhood\textquotedblright{}, under Sanskrit \emph{prátivēśa} (CDIAL 8598); \textbf{\ur{پڑوسی}}, \textquotedblleft{}a neighbour\textquotedblright{}, belongs with it.
\end{cousinweb}

\begin{grammarlens}[title={What you've built}]
One more word for getting around.
\end{grammarlens}

\subsection*{Your turn}
\begin{itemize}
\raggedright
  \item \emph{Say it:} \emph{paṛos}
  \item \emph{Say it:} \emph{paṛos}, once more
  \item \emph{Say it:} \emph{hamāre paṛos meṅ ek darzī rahtā hai.} — A tailor lives next door to us.
  \item \emph{Recall:} say the Urdu for an inhabitant, then the Urdu for the countryside, then say \emph{paṛos} again
\end{itemize}

\begin{tcolorbox}[breakable,colback=teal!4,colframe=teal!35!black,title={Before you move on}]
What does \ur{پڑوس} mean? (\textquotedblleft{}Next door; the neighbourhood\textquotedblright{}.) Say it once more.
\end{tcolorbox}
